% BHU PYQ -- Manually transcribed & error-corrected
\documentclass[11pt,a4paper]{article}

\usepackage[a4paper,top=2.5cm,bottom=2.5cm,left=2.8cm,right=2.8cm,headheight=14pt]{geometry}
\usepackage{amsmath,amssymb}
\usepackage[T1]{fontenc}
\usepackage[utf8]{inputenc}
\usepackage{lmodern}
\usepackage{microtype}
\usepackage{fancyhdr}
\usepackage{enumitem}
\usepackage{hyperref}
\usepackage{lastpage}
\usepackage{parskip}

\hypersetup{colorlinks=true,urlcolor=black,linkcolor=black,
  pdftitle={BPT-301: Optics -- BHU PYQ},pdfauthor={Banaras Hindu University}}

\pagestyle{fancy}\fancyhf{}
\renewcommand{\headrulewidth}{0.4pt}\renewcommand{\footrulewidth}{0.4pt}
\fancyhead[L]{\small   \textit{Banaras Hindu University}}
\fancyhead[R]{\small   \textit{BPT-301: Optics}}
\fancyfoot[C]{\small Page  \thepage\ of \pageref{LastPage}}


\newlist{parts}{enumerate}{2}
\setlist[parts,1]{label=(\alph*),leftmargin=*,itemsep=5pt,topsep=3pt}
\setlist[parts,2]{label=(\roman*),leftmargin=*,itemsep=3pt,topsep=2pt}

\newcommand{\pts}[1]{\hfill   \textit{[#1]}}


\begin{document}


\begin{center}
  {\large\bfseries BANARAS HINDU UNIVERSITY}\\[2pt]
  {\normalsize B.Sc. (Hons.) Semester III Examination, 2023-24}\\[6pt]
  \rule{0.8\linewidth}{0.5pt}\\[6pt]
  {\large\bfseries Paper No. BPT-301: Optics}\\[4pt]
  {\normalsize Time: 3 Hours\hfill Maximum Marks: 70}
\end{center}
\rule{\linewidth}{0.4pt}
\smallskip



\noindent   \textit{Instructions: Answer   \textbf{five} questions including Question~1 which is compulsory. Symbols have their usual meanings.}
\bigskip




\noindent   \textbf{Question 1.} Answer any   \textbf{five} of the following: \pts{5 $ \times$ 4 = 20}

\begin{parts}
  \item In a Young's double-slit experiment, interference fringes are formed using sodium light comprising two wavelengths $5890\, \text{\AA}$ and $5896\, \text{\AA}$. Obtain the regions on the screen where the fringe pattern will disappear. Assume $d = 0.5\, \text{mm}$ and $D = 100\, \text{cm}$.
  \item A plane-polarized light is incident perpendicularly on a quartz plate cut with faces parallel to the optic axis. Find the thickness of the quartz plate which introduces a phase difference of $60^\circ$ between the $e$ and $o$ rays.
  \item A left circularly polarized beam ($\lambda_0 = 5893\, \text{\AA}$) is incident normally on a calcite crystal (with its optic axis cut parallel to the surface) of thickness $0.005141\, \text{mm}$. What will be the state of polarization of the emergent beam? For calcite, $n_o = 1.65836$ and $n_e = 1.48641$.
  \item In a double-slit diffraction experiment, the slit width $b = 7  \times 10^{-3}\, \text{cm}$, the separation between identical points of the two slits $d = 7  \times 10^{-2}\, \text{cm}$, and $\lambda = 6300\, \text{\AA}$. How many interference minima will occur within the first diffraction minima on either side of the central diffraction maximum? Which orders of interference maxima will be missing?
  \item Show that a zone plate acts like a converging lens with multiple foci.
  \item How can the central fringe be seen in a Lloyd's mirror experiment? Show that the central fringe is sharper in a Lloyd's mirror than in a Fresnel biprism experiment.
  \item White light is reflected normally from a uniform oil film ($n = 1.33$). An interference maximum for $6000\, \text{\AA}$ and a minimum for $4500\, \text{\AA}$ coincide. Calculate the thickness of the film.
\end{parts}

\medskip\rule{\linewidth}{0.2pt}




\noindent   \textbf{Question 2.}

\begin{parts}
  \item In a Michelson interferometer experiment performed with a source which consists of two wavelengths $4882\, \text{\AA}$ and $4886\, \text{\AA}$, through what distance does the mirror have to be moved between two successive positions of the disappearance of the fringes? \pts{6}
  \item Discuss the intensity distribution of a Fabry-P\'erot interferometer and show that the larger the value of $R$ (reflectivity of the mirrors), the greater the difference between $I_{max}$ and $I_{min}$. In what way is it superior to the Michelson interferometer? \pts{8}
\end{parts}

\medskip\rule{\linewidth}{0.2pt}




\noindent   \textbf{Question 3.}

\begin{parts}
  \item Deduce the expression of the intensity distribution of Fraunhofer diffraction due to double slits. Discuss the nature of the intensity distribution and explain how the intensity falls off to zero on either side of the central fringe and gradually reappears. \pts{8}
  \item In a Newton's ring arrangement with $\lambda = 6.4  \times 10^{-5}\, \text{cm}$ and the radius of curvature of the convex surface is $100\, \text{cm}$, if the lens is vertically raised up, discuss the change in the ring pattern as seen through the microscope. \pts{6}
\end{parts}
\hfill   \textbf{OR}

\begin{parts}
  \item Consider a diffraction grating of width $5\, \text{cm}$ with slits of width $0.0001\, \text{cm}$ separated by a distance of $0.0002\, \text{cm}$. How many orders would be observable for $\lambda = 5.5  \times 10^{-5}\, \text{cm}$? Calculate the width of the principal maximum. Would there be any missing orders? For this diffraction grating, calculate the dispersion and resolving power in the different orders. \pts{8}
  \item What is the Rayleigh criterion of resolving power? Find the resolving power of a grating. Discuss what happens if the number of grating elements is too large. \pts{6}
\end{parts}

\medskip\rule{\linewidth}{0.2pt}




\noindent   \textbf{Question 4.}

\begin{parts}
  \item You are given a Nicol prism and a quarter-wave plate. Explain the method to produce and detect plane, elliptically, and circularly polarized light. \pts{10}
  \item What are the states of polarization when the $x$- and $y$-components of the electric field are given by the following equations:
    
\begin{parts}
      \item $E_x = E_0 \cos(\omega t + kz), \quad E_y = \frac{E_0}{2} \cos(\omega t + kz + \pi)$
      \item $E_x = E_0 \sin(\omega t + kz), \quad E_y = E_0 \cos(\omega t + kz)$
    \end{parts} \pts{4}
\end{parts}
\hfill   \textbf{OR}

\begin{parts}
  \item Explain birefringence. What are quarter- and half-wave plates? What are their uses? \pts{6}
  \item Discuss the working of a biquartz polarimeter. \pts{8}
\end{{parts}}

\medskip\rule{\linewidth}{0.2pt}




\noindent   \textbf{Question 5.}

\begin{parts}
  \item What are the differences between spontaneous and stimulated emission? Derive the relations between the Einstein coefficients. \pts{8}
  \item The length of a laser tube is $150\, \text{mm}$ and the gain factor of the laser material is $0.0005\, \text{per cm}$. If one of the cavity mirrors reflects 100 percent of the light, what is the required reflectance of the other cavity mirror? \pts{6}
\end{parts}

\vfill
\rule{\linewidth}{0.4pt}

\begin{center}\small
  \href{https://bhu-exam.vercel.app/}{Click here to visit our website}\\[4pt]
  \href{https://me-aryan.vercel.app/}{Click here to visit our contributor}\\[4pt]
  \href{https://quashan-me.vercel.app/}{Click here to visit our contributor}
\end{center}

\end{document}
